% ---------------------------------------------------------------------------
% The manuscript body, shared verbatim by the Springer (LNCS) and IEEE
% (IEEEtran) builds. Nothing class-specific may appear in this file.
%
% Each wrapper defines, before \input-ing this:
%   \wfigure / \endwfigure   a figure spanning the full text width
%   \wtable  / \endwtable    the same for tables
%   \tabscale                the font size command tables are set at
%   \para{...}               a run-in bold paragraph heading
%   \citep{...}               mapped to \cite for the LNCS/IEEEtran builds
% ---------------------------------------------------------------------------

\section{Introduction}

An operational research platform accumulates infrastructure decisions that
look, from outside the codebase, like a mistake: a machine-learning
experiment-tracking library that the code being tracked is forbidden to
import; a cache deliberately configured to ignore its own expiry policy; two
ways to fetch the same array that differ by an order of magnitude in speed
depending on which one is called for a given request shape; a data server
that answers ``not found'' for a dataset that is not, in fact, gone. Each of
these reads as a bug, a redundancy, or an oversight until the specific
measurement behind it is produced. This paper reports four such decisions
from one running system --- MarisAI, a marine ocean-data platform combining
live provider ingestion, a per-site forecasting engine, and offline
species-distribution and harmful-algal-bloom pipelines --- each with the
measurement that justified it, the mechanism that makes it necessary rather
than merely defensible, and, in three of the four, a place where the
justification itself has already partly gone stale: an upstream package's
own version pin moved, a documentation figure disagrees with the code next
to it, or a summary line's arithmetic does not say what it appears to say.

This is a short-form, practitioner-notes contribution in the SoftwareX
tradition: four self-contained technical findings under one introduction,
each independently useful to a reader who touches only one of the four
sections. What ties them together is not a shared subsystem --- one section
is about a Python dependency graph, one about a disk cache, one about a
cloud array format, one about an HTTP retry policy --- but a shared
discipline: every claim below was re-verified against a primary source this
session rather than carried forward from an existing write-up, and every
place that re-verification produced a different answer than the existing
write-up is reported rather than silently corrected, following the same
evidence-grading convention this project's other papers use.

\subsection{Contributions}

\begin{enumerate}\itemsep2pt
\item \textbf{A dependency inverted rather than added} (\S\ref{sec:mlflow}):
the backend forecasting engine is tracked by MLflow without the backend
importing MLflow, and the specific reason a lighter package is required
(a pandas/pyarrow downgrade the full package would force) is checked live
against the package's own published metadata rather than assumed --- which
shows the reason holds for the exact release pinned here, but has already
been relaxed in the package's newer releases.
\item \textbf{Two cache decisions that turn a 700-second-per-variable audit
into 0.8 seconds} (\S\ref{sec:cache}): pinning a snapshot date so a cache key
is stable, and a frozen-cache override that is correct specifically because
the window it serves cannot be revised upstream.
\item \textbf{The same cloud array format, two chunkings, opposite winners}
(\S\ref{sec:arco}): Copernicus Marine's two zarr services trade a 12-fold
speed difference depending on whether a request is one snapshot of the whole
globe or one bounded area across many timesteps, and a needed correction to
an overstated claim about how bad the wrong choice actually is.
\item \textbf{A 404 that is not proof a dataset is gone} (\S\ref{sec:erddap}):
an ERDDAP data server answers ``Currently unknown datasetID'' while
reloading a dataset it still serves, indistinguishable by status code from
a genuinely retired one, with a documented case of a real dataset believed
retired for months on that basis alone --- and a second, independent
inconsistency inside this project's own documentation about how expensive
retrying the wrong kind of failure actually is.
\end{enumerate}

\section{An MLflow Dependency Inverted, Not Added}
\label{sec:mlflow}

\para{Claim} The backend forecasting engine's training runs are tracked in
MLflow \citep{mlflow2020} without the backend ever importing MLflow. The
training script writes an immutable, plain-JSON run record to a
per-invocation directory (\texttt{\_reports/runs/<timestamp>/}, under the
forecasting engine's own model tree), using only the standard library; a
separate module on the machine-learning side,
\texttt{marine\_ml.ingest\_forecasting}, reads those directories and logs
them into MLflow after the fact. The tracked system is never aware that
tracking exists.

\para{Why the direction is inverted} The natural design --- log directly from
the training script --- would put an MLflow import inside \texttt{backend/},
which this codebase treats as a hard boundary: the API's import graph is kept
free of the modelling stack for the same reason \texttt{services/predictions.py}
serves precomputed grids rather than running inference in-process. A tracking
client imported by a training script still lands in that tree. Inverting the
dependency --- write JSON the training script already had to write; read it
from the side where MLflow is actually installed --- keeps the boundary
intact while still getting every run permanently recorded, which a fixed
filename previously did not: two training batches run back to back had
silently overwritten each other's report, the concrete failure that
motivated this design. One committed test enforces both properties directly:
it walks every Python file under \texttt{backend/} for an \texttt{import
mlflow} line, and separately checks that neither dependency file mentions
\texttt{mlflow} at all. Both checks were re-run for this paper and pass,
with \BackendMlflowOffenders{} offending imports found.

\para{Why \texttt{mlflow-skinny}, checked rather than assumed} The tracking
side pins \texttt{mlflow-skinny==\MlflowSkinnyPinned{}}, not plain
\texttt{mlflow}, alongside \texttt{pandas==\PandasPinned{}} and
\texttt{pyarrow==\PyarrowPinned{}}. The stated reason is that the full
package pins an older pandas and pyarrow ceiling that would force a
major-version downgrade underneath a multi-million-row parquet feature
store. That claim is not taken on the repository's word here: PyPI's own
release metadata for \texttt{mlflow==\MlflowMatchedVersion{}} --- the release
matching the \texttt{mlflow-skinny} version actually pinned --- was read
directly, and it does require \texttt{pandas<\MlflowMatchedPandasCeiling{}}
and \texttt{pyarrow<\MlflowMatchedPyarrowCeiling{}}, confirming the stated
reason exactly. But the same lookup against \texttt{mlflow}'s current latest
release, \MlflowLatestVersion{}, shows that ceiling has already moved:
\texttt{pandas<\MlflowLatestPandasCeiling{}} and
\texttt{pyarrow<\MlflowLatestPyarrowCeiling{}}, both of which
\PandasPinned{}/\PyarrowPinned{} already satisfy without any downgrade. The
inversion's justification is therefore correct as documented, but it is a
fact about one pinned release rather than a stable fact about ``mlflow'' ---
worth stating with a version number attached, which neither this
repository's own comments nor research.md's write-up of them originally did,
and worth re-checking at the next dependency bump rather than treated as
permanent.

\para{Tracking never fails a training run} A harmful-algal-bloom training run
takes roughly \HabRunMinutes{} minutes; losing one to a locked SQLite store or
a NaN metric would be a strictly worse outcome than an untracked run. Every
logging call in \texttt{marine\_ml.tracking} is wrapped in a context manager
that catches any exception and downgrades it to a warning, and the run-opening
code itself falls back to an inert, no-op handle if MLflow is not importable
or the tracking store cannot be reached at all --- the training pipeline
proceeds identically either way, just without a recorded run.

\para{Why it matters} A dependency that is added in the conventional
direction (import the tracker where the work happens) is the version that
gets built first and is the version that erodes an import-graph boundary a
team specifically drew for a different reason. Inverting it costs one extra
module and one regression test, and the version-checked justification above
is the kind of claim that ages --- a training script's dependency choices are
usually re-examined only when something breaks, and this one's stated reason
had already partly stopped being true one release upstream by the time this
paper checked it.

\section{Reproducing a 52-Model Audit in Under a Second per Variable}
\label{sec:cache}

\para{Claim} A companion study (\emph{Rising Skill, Falling Skill},
\texttt{research/README.md}) trains and scores \NModelsAudited{} models
(\NVariablesAudited{} variables $\times$ \NHorizonsAudited{} forecast
horizons) against two baselines, then evaluates each with a leave-one-site-out
transfer test across \NSitesAudited{} sites, for \NHeldoutFits{} held-out
fits in total (\(52 \times 24\), not a third factor folded into the 52 --- a
compressed line in the study's own README reads as one product of three
numbers when it is actually two separate multiplications, one for model
count and one for held-out fit count). The main experiment runs end to end
in \MainExperimentSeconds{}~seconds and a companion seasonality ablation in
\AblationExperimentSeconds{}~seconds, both read directly from the committed
run-metadata files this session, entirely offline.

\para{The mechanism} Two independent decisions make that possible. First,
the training window's \texttt{as\_of} snapshot date is pinned
(\PinnedAsOfDate{}), rather than computed from ``now'' as a live training run
would, so every cached provider request the audit makes has a stable cache
key across reruns; without it, each rerun would recompute a slightly
different window and re-request roughly two years of data per variable per
site from scratch. Second, an environment variable
(\texttt{MARISAI\_FROZEN\_HISTORY\_CACHE=1}) tells the disk cache to serve an
entry regardless of its normal \DiskCacheTtlHours{}-hour time-to-live. That
TTL exists because a live request asks for a window ending ``now,'' and
yesterday's answer to that question is stale by construction; a pinned,
closed historical window is the opposite case --- it cannot be revised by
the upstream provider after the fact, so re-fetching it would yield the same
bytes at real cost and occasional risk (a flapping provider can silently
return a slightly shorter series on a rerun, making two runs of the same
experiment incomparable). The override is set only by the three experiment
scripts that need it and is read nowhere on a live serving path, so a
request for ``today's'' data can never accidentally inherit a frozen answer.

\para{What is, and is not, a stored measurement} The audit's \emph{scale} ---
\NModelsAudited{} models, \NHeldoutFits{} fits, and the two total run
durations above --- is read directly from three committed report files
(\texttt{meta.json}, \texttt{loso.json}, \texttt{ablation\_meta.json}, all
under \texttt{research/data/results/}) and is reproducible by re-running the
committed pipeline. The audit's \emph{speedup}, by contrast --- from roughly
\AuditSecondsUncached{}~seconds to roughly \AuditSecondsCached{}~seconds per
variable, an \AuditSpeedupFactor{}$\times$ factor --- is stated identically
in three places (the study's README, the caching function's own docstring,
and the experiment script's module docstring) but traces to no committed
benchmark log or named timing script in any of the three. It is reported
here because it is the number that explains why a 52-model audit plus 1,248
held-out fits was practical to run at all, but it carries a different
evidence weight than the counts beside it, and is graded accordingly rather
than folded into the same claim.

\para{Why it matters} A cache that ignores its own TTL is, by default,
exactly the kind of decision a later reviewer flags as a bug: caches are
supposed to expire. The distinction that makes it correct here --- a pinned
historical window is a different object than a live one, and deserves a
different caching policy --- is not visible from the code alone; it is
visible only once the access pattern (an offline reproducibility experiment,
not a live dashboard request) is stated explicitly, which is why the
override is scoped to the three scripts that state it rather than made a
global default.

\section{Access-Pattern-Dependent Performance in ARCO Zarr Stores}
\label{sec:arco}

\para{Claim} Copernicus Marine publishes the same physics and biogeochemistry
products through two Analysis-Ready Cloud-Optimized (ARCO) zarr services ---
part of a broader shift toward cloud-native, chunked array formats for
large scientific datasets \citep{abernathey2021} --- with
opposite chunking layouts: \texttt{arco-geo-series} stores huge latitude/
longitude chunks with one timestep each, and \texttt{arco-time-series} stores
fine latitude/longitude chunks with huge time chunks. Which one is faster
depends entirely on the shape of the request, and this one codebase makes
requests of both shapes: a whole-globe, single-timestep snapshot for map-tile
rendering and the global habitat-model build, and a bounded-area,
many-timestep series for the regional downloader and the per-site
forecasting engine. Using the wrong service for a given request shape is a
correctness-preserving but severe performance regression, not a crash, which
is exactly the kind of defect that survives unnoticed until someone times it.

\para{The one direction actually measured} \texttt{docs/ml-notes.md} records
a direct comparison, dated \ArcoMeasuredDate{}, for one year of global
monthly physics coarsened to 0.25\si{\degree}: \texttt{arco-geo-series}
completes in \ArcoGeoSeriesSeconds{}~seconds against
\texttt{arco-time-series}'s \ArcoTimeSeriesSeconds{}~seconds for the
identical request --- an \ArcoSlowdownFactor{}$\times$ difference for
picking the service built for the opposite access pattern. The reversal for
the opposite request shape (a bounded area, many timesteps, where
\texttt{arco-time-series}'s fine spatial chunking and coarse time chunking
should win) is not a separately timed measurement in any source read for
this paper; it is stated as a direct, mechanistic consequence of the two
services' documented, opposite chunk layouts, and is treated here as
asserted rather than measured. Three access patterns exist in this one
codebase (map-tile snapshot, regional download, whole-globe forecast grid),
and the service is selected per call rather than as a module-wide default
specifically so a future caller cannot silently inherit the wrong one for
its own request shape.

\para{A second, independent hazard: a missing depth bound} The
\ReanalysisDepthLevels{}-level \texttt{\_my\_} reanalysis products carry
every one of those levels even when only the surface is needed. Passing
\texttt{minimum\_depth}/\texttt{maximum\_depth} into \texttt{open\_dataset}
so the server subsets by depth as well as by space and time turns a request
that otherwise did not finish in \ReanalysisUnboundedMinutes{} minutes into
one that completes in roughly \ReanalysisBoundedSeconds{} seconds. This is a
different mechanism from the service-selection finding above --- a missing
server-side filter, not a choice between two services with different
chunking --- and is kept as a separate claim rather than merged into the
same headline number.

\para{A claim this paper does not repeat} Both this project's own working
notes (\texttt{CLAUDE.md}) and research.md's write-up of them describe
picking the wrong ARCO service as turning ``a $\sim$10-second fetch into
minutes/hours'' (\texttt{CLAUDE.md}) or ``hours'' outright (research.md).
Neither figure is supported by any measurement located this session: the one
timed comparison above is \ArcoGeoSeriesSeconds{}s against
\ArcoTimeSeriesSeconds{}s, an \ArcoSlowdownFactor{}$\times$ factor and not
hours, and the depth-bound hazard's worst measured case is stated as ``did
not finish in \ReanalysisUnboundedMinutes{} minutes,'' which bounds how long
an engineer waited before giving up, not how long the request would actually
have taken. This paper reports only the two measured-or-explicitly-stated
pairs above rather than the unquantified ``hours'' framing.

\para{Why it matters} A cloud-native array format advertised as one product
can hide two internally inconsistent access patterns behind one API surface,
and the two can differ by an order of magnitude for the same logical
request. The fix costs one keyword argument per call site; the hazard is
that the module-level default silently applies to every future caller unless
the choice is forced at each call site, which is why this codebase treats it
that way.

\section{ERDDAP Answers 404 While Reloading a Dataset}
\label{sec:erddap}

\para{The corrected belief} For months, NOAA CoastWatch's ERDDAP-served
\texttt{GEBCO\_2020} bathymetry dataset was treated in this codebase as
permanently retired, because every request for it returned HTTP 404. That
belief was wrong. Re-probed \GebcoReprobeDate{}, \texttt{GEBCO\_2020} returns
200 with correct depths. What actually happens, confirmed by the fix
committed that same day, is that the whole server flaps: over one probing
session it went fully 503, then served \texttt{GEBCO\_2020} at 200 for
minutes, then answered 404 for it again for roughly \ErddapFlapSeconds{}
seconds before recovering. ERDDAP unloads a dataset while reloading it and
answers 404 ``Currently unknown datasetID'' during that window ---
indistinguishably, by status code and message, from a dataset that has
actually been removed. The production bathymetry provider now reads
\texttt{gebco2021} from a different host (Ifremer's ERDDAP), chosen because
it is the newer GEBCO release and happened to stay up throughout the
flapping session, not because CoastWatch's copy is actually gone.

\para{The retry policy this forces} A 404 cannot be treated as either always
permanent or always transient without cost on one side or the other, so the
classification is asymmetric and specific to the status code. Confirmed
directly against \texttt{backend/forecasting/history.py}'s retry
classifier and its currently-passing test suite: a \texttt{404} is retried
exactly \RetryAttemptsFourZeroFour{} times, with
\RetrySecondsFourZeroFour{}~seconds of backoff, because a request landing
inside a reload window recovers on the second attempt; a \texttt{400} (and
the other permanent 4xx codes, 403/410/422) gets \RetryAttemptsFourHundred{}
attempt and no retry at all, because those mean the request itself is
malformed and will stay malformed; and a \texttt{503} (with the other
transient codes, 408/425/429, and any 5xx) is retried up to
\RetryAttemptsFiveZeroThree{} times with \RetrySecondsFiveZeroThree{} seconds
of total backoff, because a fully overloaded or restarting server is exactly
the case unlimited-within-budget retrying exists for.

\para{Both costs, measured} Getting this wrong costs real time in both
directions, and both directions are documented with a specific number rather
than left as a general concern. Retrying a permanent failure with the old,
uniform three-attempt policy added the full \RetryPenaltySeconds{}-second
backoff (\RetryBackoffFirstSeconds{}s + \RetryBackoffSecondSeconds{}s) to
\emph{every} request against that dataset. The two sources closest to the
implementation --- \texttt{history.py}'s own code comment and the regression
test built specifically around this scenario --- both state this turned a
warm \WarmForecastCodeSeconds{}-second response into
\WarmForecastAfterBadRetrySeconds{} seconds; this project's own documentation
page and research.md's write-up of it instead state a warm
\WarmForecastDocsSeconds{}-second response becoming the same
\WarmForecastAfterBadRetrySeconds{} seconds. Both pair with the identical
\RetryPenaltySeconds{}-second penalty, and only one addition is close to the
stated total: \(\WarmForecastCodeSeconds{} + \RetryPenaltySeconds{} = 14\),
against \(\WarmForecastDocsSeconds{} + \RetryPenaltySeconds{} = 10.2\) ---
neither reaches \WarmForecastAfterBadRetrySeconds{} exactly, but the
code-and-test figure is the one the arithmetic actually supports, and it is
reported here as primary for that reason, with both figures stated rather
than the disagreement silently resolved. In the other direction, retrying
too conservatively costs geography rather than latency: a single transient
ERDDAP 503 during one training run cost \ErddapCostAffectedPoints{} of
\ErddapCostTotalPoints{} training points --- a quarter of that model's
spatial coverage, including the entire Gulf Stream and Coral Sea --- from a
provider that supplies one constant column and had recovered again within
seconds.

\para{A live loose end} One committed test's own prose has not caught up
with the correction above. \texttt{backend/tests/test\_download\_gebco.py}
still narrates \texttt{GEBCO\_2020} as a dataset ``the server retired'' that
``is gone,'' language describing the pre-\GebcoReprobeDate{} belief rather
than the flapping-server explanation recorded in the provider module's own
comment. The test itself is not wrong --- it still correctly asserts the
provider requests Ifremer's \texttt{gebco2021} rather than
\texttt{GEBCO\_2020}, which remains true for the unrelated reason that
Ifremer stayed available throughout the flap --- but its docstring was never
updated to match the correction shipped elsewhere in the same area of the
codebase, a small, live instance of exactly the documentation drift this
section's other paragraph reports.

\para{Why it matters} A 404 is conventionally treated as the internet's most
unambiguous status code: the thing you asked for does not exist. On an
ERDDAP host mid-reload it means something closer to ``ask again in a
moment,'' and the two cannot be told apart from the response alone. Getting
the retry policy wrong in either direction has a real, opposite cost --- a
slow page on one side, missing training geography on the other --- and both
costs were only visible once someone measured them rather than picked a
retry count that felt reasonable.

\begin{wtable}
\centering
\tabscale
\caption{The four findings, what each decision looks like without its
measurement, and this paper's evidence grade for the headline numbers (see
\texttt{shared/generated/macros.tex} for the full per-group grading).
Artifact = a committed report/data file or a currently-passing test
reproduces the number on demand; Documentation = descriptive, reproducible
prose with no stored benchmark output; Dated = a single, named-date
measurement with no committed artifact; External = checked live against a
third party's own published metadata.}
\label{tab:findings}
\begin{tabular}{p{0.04\linewidth}p{0.22\linewidth}p{0.47\linewidth}p{0.17\linewidth}}
\toprule
\S & Looks like & Key numbers & Evidence \\
\midrule
4.1 & an unneeded extra package & \MlflowMatchedVersion{} needs pandas
$<$\MlflowMatchedPandasCeiling{}, pyarrow $<$\MlflowMatchedPyarrowCeiling{};
\MlflowLatestVersion{} relaxes this to $<$\MlflowLatestPandasCeiling{}/
$<$\MlflowLatestPyarrowCeiling{} & Artifact, External \\
4.2 & a cache ignoring its own expiry & \NModelsAudited{} models,
\NHeldoutFits{} fits; \AuditSecondsUncached{}s to \AuditSecondsCached{}s
per variable & Artifact (counts), Documentation (speedup) \\
4.3 & an arbitrary service choice & \ArcoGeoSeriesSeconds{}s vs
\ArcoTimeSeriesSeconds{}s, one direction timed; depth bound
\ReanalysisUnboundedMinutes{}min to \ReanalysisBoundedSeconds{}s & Dated \\
4.4 & an over-cautious retry count & 404: \RetryAttemptsFourZeroFour{}
attempts/\RetrySecondsFourZeroFour{}s; 400: \RetryAttemptsFourHundred{};
503: \RetryAttemptsFiveZeroThree{}/\RetrySecondsFiveZeroThree{}s;
\WarmForecastCodeSeconds{}s to \WarmForecastAfterBadRetrySeconds{}s
misrouted & Artifact (policy), Documentation (cost) \\
\bottomrule
\end{tabular}
\end{wtable}

\section{Discussion}

Read together, the four findings share a shape that is easy to miss reading
them one at a time: in every case, the decision that looks wrong from the
code alone is the one a specific, dated measurement made necessary, and in
three of the four, some part of the write-up describing that measurement has
already drifted from the primary source by the time this paper re-checked
it. The MLflow inversion's stated justification is exactly correct for the
release pinned, and already one release out of date for the package in
general. The 52-model audit's headline speedup is real but is prose, not an
artifact, sitting next to counts that are. Two independent sources inside
this same codebase --- a source-code comment and its own regression test on
one side, a documentation page and research.md's copy of it on the other ---
disagree by a factor of nearly three on how expensive a bad retry policy
actually was, and the arithmetic embedded in the claim itself is what
settles which one to believe rather than either source's authority. None of
these is the kind of drift that changes the underlying finding; each is the
kind that changes a specific number a reader might otherwise quote with more
confidence than it has earned.

That pattern is not unique to infrastructure code, and this project's
adjacent work on silent, plausible failure modes in this same platform's
modelling pipelines found the identical structure at the modelling layer: a
number that is real, computed correctly, and licensed by convention or by a
prior write-up to mean more than a second check shows it to mean. The
version-checked dependency claim, the counts-versus-prose distinction in the
cache speedup, and the 6-versus-2.2-second reconciliation in
\S\ref{sec:erddap} are three small instances of exactly that discipline
applied to this paper's own sources rather than to the system being
described --- which is the point: a paper reporting that infrastructure
decisions need re-verification against their primary sources should not
itself rest on an unverified secondary one.

\section{Conclusion}

Four infrastructure decisions from one operational marine-data platform ---
an inverted MLflow dependency, a cache that deliberately ignores its own
expiry, a per-call choice between two zarr services with opposite chunking,
and an asymmetric retry policy for one ambiguous HTTP status code --- each
looks like a defect or an arbitrary choice until the measurement behind it
is produced, and each measurement is reported here with the evidence grade
it actually earns rather than a uniform confidence. Three of the four
carried a documentation figure that had already drifted from its own primary
source by the time this paper checked it: an upstream package's version pin
that has since been relaxed, a per-variable speedup that is prose rather
than a stored benchmark, and a retry-cost example that disagrees with itself
by nearly a factor of three depending on which committed file is read. None
of the four findings changes as a result; every one of the three corrections
is reported explicitly rather than folded quietly into a rounder number,
because a paper about verifying infrastructure decisions against their
sources has to hold its own numbers to the same standard.
